\tlatex
\begin{lcom}{0}%
\begin{cpar}{0}{F}{F}{0}{0}{}%
 This spec uses the \ensuremath{brace} \ensuremath{conjunctions} directive:
 each definition whose
 body is a bulleted conjunction list is typeset like a piecewise function,
 with a tall left brace in place of the \ensuremath{\.{\land}} bullets and
 the defined name and
 its \ensuremath{\.{\defeq}} centered vertically on the list. The last
 definitions show the
 cases that are typeset as usual.
\end{cpar}%
\end{lcom}%
\@x{}\moduleLeftDash\@xx{ {\MODULE} ConjunctionBraces}\moduleRightDash\@xx{}%
\@x{ {\EXTENDS} Naturals ,\, FiniteSets}%
\@pvspace{8.0pt}%
\@x{ {\CONSTANTS} Proc ,\, MaxCount}%
\@x{ {\VARIABLES} pc ,\, count ,\, msgs}%
\@pvspace{8.0pt}%
\@x{}%
\@y{\@s{0}%
 The list follows the \ensuremath{\.{\defeq}}.
}%
\@xx{}%
\@bx{ TypeOK \.{\defeq}}{%
 \@x{\phantom{ \.{\land}} pc \.{\in} [ Proc \.{\rightarrow} \{\@w{idle}
 ,\,\@w{busy} ,\,\@w{done} \} ]}%
 \@x{\phantom{ \.{\land}} count \.{\in} [ Proc \.{\rightarrow} 0 \.{\dotdot}
 MaxCount ]}%
 \@x{\phantom{ \.{\land}} msgs \.{\subseteq} Proc \.{\times} ( 0 \.{\dotdot}
 MaxCount )}%
}%
\@pvspace{8.0pt}%
\@x{}%
\@y{\@s{0}%
 The list starts on the next line; it moves up beside the
 \ensuremath{\.{\defeq}}.
}%
\@xx{}%
\@bx[0]{ Init \.{\defeq}}{%
\@x{\phantom{ \.{\land}} pc \.{=} [ p \.{\in} Proc \.{\mapsto}\@w{idle} ]}%
\@x{\phantom{ \.{\land}} count \.{=} [ p \.{\in} Proc \.{\mapsto} 0 ]}%
\@x{\phantom{ \.{\land}} msgs \.{=} \{ \}}%
}%
\@pvspace{8.0pt}%
\@x{}%
\@y{\@s{0}%
 Conjuncts may span lines, contain disjunction lists and end-of-line
}%
\@xx{}%
\@x{}%
\@y{\@s{0}%
 comments, and a nested definition\mbox{'}s list gets its own brace.
}%
\@xx{}%
\@bx[0]{ Step ( p ) \.{\defeq}}{%
\@x{\phantom{ \.{\land}} pc [ p ] \.{=}\@w{idle}}%
\@x{\phantom{ \.{\land}} \.{\lor} count [ p ] \.{<} MaxCount}%
\@y{\@s{0}%
 room for another message
}%
\@xx{}%
\@x{ \.{\lor} msgs \.{=} \{ \}}%
\@bx{\phantom{ \.{\land}} \.{\LET} next \.{\defeq}}{%
 \@x{\phantom{ \.{\land}} count \.{'} \.{=} [ count {\EXCEPT} {\bang} [ p ]
 \.{=} @ \.{+} 1 ]}%
 \@x{\phantom{ \.{\land}} msgs \.{'} \.{=} msgs \.{\cup} \{ {\langle} p ,\,
 count [ p ] {\rangle} \}}%
}%
\@x{ \.{\IN} next}%
 \@x{\phantom{ \.{\land}} pc \.{'} \.{=} [ pc {\EXCEPT} {\bang} [ p ]
 \.{=}\@w{busy} ]}%
}%
\@pvspace{8.0pt}%
\@x{}%
\@y{\@s{0}%
 An invariant with one wide conjunct: it does not fit beside the long
}%
\@xx{}%
\@x{}%
\@y{\@s{0}%
 name, so the name is on its own line and the brace below it.
}%
\@xx{}%
\@bx[0]{ EveryProcessThatIsDoneHasSentAtLeastOneMessage \.{\defeq}}{%
 \@x{\phantom{ \.{\land}} \A\, p \.{\in} Proc \.{:} pc [ p ] \.{=}\@w{done}
 \.{\implies} \E\, m \.{\in} msgs \.{:} m [ 1 ] \.{=} p \.{\land} m [ 2 ]
 \.{<} MaxCount}%
 \@x{\phantom{ \.{\land}} Cardinality ( msgs ) \.{\leq} Cardinality ( Proc )
 \.{*} ( MaxCount \.{+} 1 )}%
}%
\@pvspace{8.0pt}%
\@x{}%
\@y{\@s{0}%
 Typeset as usual: a single conjunct, a comment between the
 \ensuremath{\.{\defeq}} and the
}%
\@xx{}%
\@x{}%
\@y{\@s{0}%
 list, and a body that is not just the list.
}%
\@xx{}%
\@x{ Single \.{\defeq} \.{\land} msgs \.{=} \{ \}}%
\@x{ Commented \.{\defeq}}%
\@y{\@s{0}%
 every process is idle
}%
\@xx{}%
\@x{ \.{\land} pc \.{=} [ p \.{\in} Proc \.{\mapsto}\@w{idle} ]}%
\@x{ \.{\land} msgs \.{=} \{ \}}%
 \@x{ Quantified \.{\defeq} \A\, p \.{\in} Proc \.{:} \.{\land} count [ p ]
 \.{\leq} MaxCount}%
\@x{ \.{\land} pc [ p ] \.{\in} \{\@w{idle} ,\,\@w{busy} ,\,\@w{done} \}}%
\@x{}\bottombar\@xx{}%
\@pvspace{8.0pt}%
